\documentclass[11pt,twoside]{article}

\usepackage[margin=1in]{geometry}
\usepackage{amsmath,amssymb,amsthm,mathtools}
\usepackage[inline,shortlabels]{enumitem}
\usepackage{xcolor}
\usepackage{microtype}

% Revision date: maintained by scripts/build_lecture_notes.py. Keep this command.
\newcommand{\lecturerevised}{not recorded}
\usepackage{eso-pic}
\AddToShipoutPictureBG{%
  \AtPageLowerLeft{%
    \put(\LenToUnit{0.5\paperwidth},\LenToUnit{0.5in}){%
      \makebox(0,0){\normalfont\footnotesize\color{black!60}Last revised: \lecturerevised}%
    }%
  }%
}

\usepackage[square,authoryear]{natbib}
\usepackage[colorlinks=true,linkcolor=blue!60!black,
  urlcolor=blue!60!black,citecolor=blue!60!black]{hyperref}
\usepackage[nameinlink,capitalize]{cleveref}
\crefname{exercise}{Exercise}{Exercises}
\Crefname{exercise}{Exercise}{Exercises}

% When including shared exercise chunks in the repository, uncomment:
% \usepackage{../exercises/course-exercises}

% Short refinements that may be skipped on a first reading.
\usepackage{enotez}
\DeclareInstance{enotez-list}{lecture}{list}{format=\normalsize,number={#1.}}
\setenotez{list-name=Endnotes,list-style=lecture,backref=true}
\crefname{endnote}{Endnote}{Endnotes}
\Crefname{endnote}{Endnote}{Endnotes}


\definecolor{courseblue}{HTML}{17365D}
\newcommand{\E}{\mathbb E}
\newcommand{\I}{\mathbb I}
\renewcommand{\P}{\mathbb P}
\newcommand{\R}{\mathbb R}
\newcommand{\KL}{D}
\newcommand{\TV}{\mathrm{TV}}
\makeatletter
\@for\@letter:=A,B,C,D,E,F,G,H,I,J,K,L,M,N,O,P,Q,R,S,T,U,V,W,X,Y,Z\do{%
  \expandafter\edef\csname c\@letter\endcsname{\noexpand\mathcal{\@letter}}}
\makeatother
\newcommand{\lecturenumber}{N}

\newtheorem{theorem}{Theorem}[section]
\newtheorem{proposition}[theorem]{Proposition}
\newtheorem{corollary}[theorem]{Corollary}
\theoremstyle{definition}
\newtheorem{definition}[theorem]{Definition}
\newtheorem{example}[theorem]{Example}
\theoremstyle{remark}
\newtheorem{remark}[theorem]{Remark}

\renewcommand{\thepage}{\lecturenumber--\arabic{page}}
\renewcommand{\thesection}{\lecturenumber.\arabic{section}}
\renewcommand{\theequation}{\lecturenumber.\arabic{equation}}

\pagestyle{myheadings}
\markboth{Lecture \lecturenumber: Lecture title}
{Lecture \lecturenumber: Lecture title}

% Course-wide notation convention:
% Every vector is a column vector. If a matrix stores one vector per position,
% row i contains the transpose of that vector. Write every row-shaped vector
% explicitly as the transpose of a column vector.
% Use n for the number of independent training examples. Use sentence case in
% optional theorem titles. Use example environments for worked examples that
% are referenced later. Reserve \boxed{} for literal boxes in diagrams.

\begin{document}

\begin{center}
\fcolorbox{courseblue}{white}{%
\begin{minipage}{0.92\textwidth}
\vspace{1mm}
\textbf{CMPUT 654: Mathematical Foundations of Modern AI Systems}
\hfill Fall 2026

\vspace{3mm}
\begin{center}
{\Large\bfseries Lecture \lecturenumber: Lecture title}

\vspace{1mm}
Lecture date
\end{center}

\vspace{2mm}
\textit{Lecturer: Csaba Szepesv\'ari}
\hfill
\textit{Scribe: Student name}
\vspace{1mm}
\end{minipage}}
\end{center}

% Follow lecture-writing-standards.md from the first draft onward.
% Organize the body around the lecture's mathematical development. Replace
% the sample headings below with headings appropriate to the lecture.
% Let equations carry transparent algebraic operations. Use prose to explain
% purpose, structural meaning, design choices, and consequences.

\noindent\textbf{Reading guide.}
State the lecture's purpose and prerequisites, identify the material for a
first reading, and point to optional material. When relevant, distinguish
what was covered in class from subsequent additions.

\section{Main question and setup}

State the question developed in the lecture. Define every mathematical object
before using it. For each matrix or tensor, state what every axis indexes.

\section{Central result or calculation}

State the main result or calculation precisely and reconstruct its argument.
Use additional sections when the lecture contains several substantial blocks.

\section{Interpretation}

Explain what the mathematics says about the AI systems discussed in the
lecture.

% A separate Optional reading section, if needed, goes here.

\section{Take-home messages}
\label{sec:take-home}

\begin{itemize}[leftmargin=*]
\item Record the main mathematical and conceptual conclusions compactly.
\end{itemize}

% From the first draft, use endnotes for short refinements that may be skipped
% on a first reading. Keep necessary definitions and assumptions in the main text;
% put substantial derivations in an appendix. Add a note beside its claim:
% A statement.\endnote{A short refinement.}
% This command prints nothing if no endnotes were defined.
\clearpage
\printendnotes

\section*{Bibliographical notes}

Explain the origin or significance of each cited reference and its connection
to the material developed above.

\section*{Glossary}

% Include a term when the note gives it a technical meaning needed to follow
% the lecture or when the term is likely to recur later in the course. Include
% persistent notation and diagram conventions. Ordinary mathematical words,
% implementation details, named models, libraries, and terms confined to
% optional background can be omitted. Define every term at first use in the
% main text; the glossary is a recap.

\begin{description}[style=nextline,leftmargin=1.5em,labelindent=0pt]
\item[Term]
Definition in the sense used in the lecture.
\end{description}

\section{Exercises}
\label{sec:exercises}

Collect focused exercises in lecture order, splitting long sets when they
develop distinct ideas. Put hints after the questions. Keep matching model
solutions in the private exercise files.

% Load course-exercises above, then replace this placeholder with:
% \exerciseinstructions
% \setcounter{courseexercise}{0}
% \input{../exercises/<exercise-id>.tex}
% Add each exercise to ../exercises/catalog.json in the same order.

% Substantial deferred proofs and background go here, when needed:
% \appendix
% \section{Title of the deferred material}
% Restate each result immediately before its proof environment. Order the
% appendices by first substantive use in the main text. Exercise solutions
% belong in private solution files.

% Put the ordinary reference list at the very end, after any appendices.
\clearpage
\begin{thebibliography}{99}
\bibitem[Author(Year)]{reference-key}
Author.
\newblock Title.
\newblock Venue, year.
\newblock \url{https://example.com}.
\end{thebibliography}

\end{document}
